% Pacote LaTeX — Relatório Final do Projeto Cripto Agent (Protocolo Chave-Organização v3)
% Compilar com: pdflatex relatorio.tex (duas vezes para referências)
% Ou: make
\documentclass[12pt,a4paper]{article}

\usepackage[utf8]{inputenc}
\usepackage[T1]{fontenc}
\usepackage[brazil]{babel}
\usepackage{geometry}
\geometry{top=2.5cm, bottom=2.5cm, left=2.5cm, right=2.5cm}

\usepackage{xcolor}
\usepackage{tikz}
\usetikzlibrary{arrows.meta, positioning, shapes.geometric, calc, fit, backgrounds}

\usepackage{booktabs}
\usepackage{array}
\usepackage{colortbl}
\usepackage{amsmath, amssymb, amsthm}
\usepackage{enumitem}
\usepackage{fancyhdr}
\usepackage{hyperref}
\usepackage{listings}
\usepackage{caption}

% ---------------- Cores ----------------
\definecolor{corTitulo}{HTML}{1A3A5C}
\definecolor{corSecao}{HTML}{2C3E50}
\definecolor{corTexto}{HTML}{34495E}
\definecolor{corAgent}{HTML}{E8F4FD}
\definecolor{corEscrow}{HTML}{FDF2E8}
\definecolor{corKey}{HTML}{FFFFFF}
\definecolor{corBorder}{HTML}{2C3E50}
\definecolor{corVerde}{HTML}{27AE60}
\definecolor{corAzul}{HTML}{2980B9}
\definecolor{corRoxo}{HTML}{8E44AD}
\definecolor{corLaranja}{HTML}{E67E22}
\definecolor{corVermelho}{HTML}{C0392B}
\definecolor{corDestaque}{HTML}{EEF2FF}

% ---------------- Listings ----------------
\lstset{
  language=Python,
  basicstyle=\ttfamily\small,
  frame=single,
  breaklines=true,
  keywordstyle=\color{corSecao}\bfseries,
  stringstyle=\color{corVerde},
  commentstyle=\color{gray!60},
  numbers=left,
  numberstyle=\tiny\color{gray},
  showstringspaces=false
}

% ---------------- Cabeçalho / Rodapé ----------------
\pagestyle{fancy}
\fancyhf{}
\renewcommand{\headrulewidth}{0.4pt}
\fancyhead[C]{\small Cripto Agent — Relatório Final do Projeto}
\fancyfoot[C]{\small Página \thepage}

% ---------------- Estilo de seções ----------------
\newcommand{\sect}[1]{
  \vspace{0.6cm}
  {\color{corTitulo}\Large\bfseries #1}
  \vspace{0.2cm}
  {\color{corBorder}\hrule height 0.8pt}
  \vspace{0.3cm}
}

\newcommand{\subsect}[1]{
  \vspace{0.4cm}
  {\color{corSecao}\large\bfseries #1}
  \vspace{0.15cm}
}

\newcommand{\capsula}[1]{{\color{corTexto}\small #1}}

\begin{document}

% ================= CAPA =================
\begin{titlepage}
\begin{center}
{\color{corTitulo}\Huge\bfseries RELATÓRIO FINAL DO PROJETO}\\[1.2cm]

{\color{corSecao}\huge\bfseries Cripto Agent}\\[0.4cm]
{\large Protocolo de Comunicação Multiagente com Escrow Organizacional}\\[0.6cm]

{\color{corTexto}\large Protocolo Chave-Organização v3 — Criptografia de Mensagem com Chave de Descriptografia Interna da Empresa}\\[1.2cm]

{\small Construído com base no template: \emph{Modelo Estrutural para Relatório Final do Projeto}\\[0.2cm]
Versão: rascunho (draft) — para homologação e revisão}\\[1.4cm]

{\color{corTexto}\normalsize
Autor: I.S.A.A.C. (Inteligência Autônoma de Suporte e Análise Computacional)\\
Criador: Matheus Soranço\\
Data: \today
}
\end{center}
\end{titlepage}

% ================= 1. NOME =================
\sect{1. Nome do Projeto}

\textbf{Cripto Agent — Protocolo Chave-Organização v3.}

Um protocolo de comunicação de interoperabilidade multiagente para ambientes organizacionais, capaz de comunicar agente-a-agente, com suporte a tools e skills, identidade independente de máquina ou perfil do sistema operacional, e criptografia de mensagem porque o tráfego navega em rede organizacional.

% ================= 2. BREVE DESCRIÇÃO =================
\sect{2. Breve Descrição}

\subsubsection*{a) Função principal}
O programa permite que agentes de software, distribuídos entre diferentes máquinas, serviços e contas de usuário de uma organização, descubram capacidades e solicitem a execução de ferramentas (\emph{tools}) e habilidades (\emph{skills}) de outros agentes, com autenticação, controle de acesso e criptografia durante a transmissão. A identidade do agente não depende do computador nem do perfil do sistema operacional; é estabelecida por chaves criptográficas cadastradas administrativamente.

\subsubsection*{b) Funções específicas relevantes}
\begin{itemize}[leftmargin=*]
  \item Consulta de catálogo de capacidades autorizadas por remetente.
  \item Invocação de tools (operações concretas) e skills (capacidades de nível mais alto).
  \item Verificação de permissão (negar por padrão) antes de executar qualquer capacidade.
  \item Validação de destinatário, organização, remetente, versão do protocolo e prazo de validade da mensagem.
  \item Registro persistente de identificadores de mensagem para rejeição de replay (SQLite).
  \item Retorno de resultados e erros estruturados em JSON.
  \item Criptografia de mensagem ponta-a-ponta com escrow organizacional: a empresa possui uma chave de descriptografia para segurança interna / auditoria.
\end{itemize}

\subsubsection*{c) Usuários atendidos primordialmente}
Desenvolvedores de sistemas multiagente, integradores de serviços e administradores de infraestrutura de organizações que implantam agentes autônomos. O programa é direcionado a perfis técnicos capazes de gerenciar chaves criptográficas e políticas de autorização.

\subsubsection*{d) Natureza do programa}
Prova de conceito (protótipo funcional) de um protocolo de comunicação original, com estrutura criptográfica própria (estilo chaves). Não é um serviço pronto para produção; requer homologação entre máquinas reais e recursos operacionais antes de adoção em processos organizacionais reais.

\subsubsection*{e) Ressalvas}
\begin{itemize}[leftmargin=*]
  \item Ainda não há homologação entre máquinas distintas em rede organizacional real.
  \item O contrato de comunicação é próprio (versão 3); adaptadores A2A/MCP ainda não foram implementados.
  \item A criptografia protege o conteúdo da mensagem; não elimina riscos de comprometimento de endpoints, roubo de chaves privadas ou uso indevido de uma capacidade autorizada.
  \item Chaves privadas não estão protegidas por senha nesta versão; em produção a chave de escrow da empresa deve ser guardada em HSM com ACLs rigorosas.
\end{itemize}

% ================= 3. VISÃO DE PROJETO =================
\sect{3. Visão de Projeto — Cenários}

Esta seção apresenta quatro cenários (dois positivos e dois negativos) que orientam o projeto do software, sua interpretação pelo usuário e sua evolução futura, conforme o modelo estrutural do relatório.

\subsubsection*{Cenário Positivo 1 — Consulta e execução de uma tool}
Uma desenvolvedora precisa integrar dois agentes de serviços internos. O administrador cadastra as identidades dos agentes e autoriza o primeiro a usar a \texttt{echo} do segundo. A desenvolvedora configura os endereços e inicia a comunicação. O agente solicitante consulta o catálogo do destinatário e recebe apenas as capacidades às quais possui acesso. Em seguida, envia um texto para a \texttt{echo}. O destinatário verifica a assinatura Ed25519 do remetente, a organização, o destino, o prazo de validade e a permissão para executar a tool. Após as verificações, processa a solicitação e devolve o texto pela mensagem cifrada.

Este cenário evoca o fluxo de descoberta e invocação de capacidade com autenticação e autorização explícitas, funcionando como bússola para o projeto do protocolo.

\subsubsection*{Cenário Positivo 2 — Delegação com auditoria da empresa (escrow)}
Um integrador precisa que um agente delegue a análise de um texto a outro serviço, e a organização exige que qualquer mensagem sensível possa ser revista internamente para fins de segurança. O agente solicitante cifra a mensagem com o protocolo v3: o payload é protegido com AES-256-GCM e a chave simétrica é encapsulada tanto para o destinatário quanto para a chave de escrow da organização. O destinatário executa a \texttt{text.stats} e retorna a quantidade de palavras e caracteres. A empresa, utilizando sua chave privada de escrow, pode descriptografar a mesma mensagem independentemente, sem precisar da chave privada do destinatário, e auditar o conteúdo para segurança interna.

\subsubsection*{Cenário Negativo 1 — Capacidade necessária indisponível}
Uma equipe tenta integrar um agente que precisa de uma habilidade ainda não registrada no destinatário. Ao consultar o catálogo, não encontra a capacidade desejada e não consegue concluir sua tarefa. O protocolo não instala automaticamente novas habilidades nem transfere código executável entre agentes. A equipe deve implementar a capacidade no serviço responsável, definir seu contrato e solicitar as permissões apropriadas.

\subsubsection*{Cenário Negativo 2 — Perda da resposta e risco de replay}
Um agente envia uma solicitação, mas perde a conexão antes de receber a resposta. Ao reenviar a mesma mensagem, o destinatário rejeita como repetição, pois o identificador já foi registrado. Como a versão atual não mantém um histórico consultável de resultados, o solicitante pode não saber se a primeira operação terminou. Para operações com efeitos externos, uma nova tentativa exige verificar o estado do sistema antes de gerar outro identificador.

% ================= 4. DOCUMENTAÇÃO TÉCNICA =================
\sect{4. Documentação Técnica do Projeto}

\subsubsection*{4.1 Requisitos funcionais e não-funcionais}

\begin{table}[h!]
\centering
\small
\begin{tabular}{p{1.6cm}p{13cm}}
\toprule
\rowcolor{corTitulo}\textcolor{white}{\textbf{ID}} & \textcolor{white}{\textbf{Requisito}} \\
\midrule
RF01 & Autenticar os agentes participantes por chaves criptográficas (Ed25519 identidade / X25519 acordo). \\
RF02 & Associar chaves públicas cadastradas às identidades dos agentes (cadastro administrado). \\
RF03 & Disponibilizar catálogo de capacidades filtrado por permissão de cada solicitante. \\
RF04 & Permitir a invocação de tools e skills registradas localmente, sem carregar código remoto. \\
RF05 & Verificar permissão (remetente, tipo, capacidade) antes de executar; negar por padrão. \\
RF06 & Validar organização, remetente e destinatário da mensagem contra o cabeçalho assinado. \\
RF07 & Rejeitar mensagens vencidas, incompatíveis com a versão do protocolo ou fora da janela temporal. \\
RF08 & Registrar identificadores de mensagem e rejeitar mensagens repetidas durante sua validade (persistência SQLite). \\
RF09 & Retornar resultados e erros estruturados em JSON. \\
RF10 & Permitir que a organização descriptografe qualquer mensagem via chave privada de escrow (segurança interna). \\
\bottomrule
\end{tabular}
\caption{Requisitos funcionais do protocolo v3.}
\end{table}

\textbf{Requisitos não funcionais:}
\begin{itemize}[leftmargin=*]
  \item Transporte independente: o envelope é um JSON autocontido (pode trafegar por HTTP/HTTPS, filas ou qualquer canal de rede).
  \item Confidencialidade: AES-256-GCM com chave efêmera por mensagem; apenas destinatário e empresa recuperam a chave.
  \item Integridade e autenticidade: assinatura Ed25519 sobre o cabeçalho completo (inclui ambos os encapsulamentos de chave).
  \item Controle de acesso: negação por padrão; capacidade só executa com permissão explícita.
  \item Limite de corpo: 65.536 bytes por mensagem. Validade curta e UUID para controle de replay.
  \item Identidade independente de máquina e perfil do Windows: a identidade é a chave pública cadastrada.
\end{itemize}

\subsubsection*{4.2 Arquitetura e modelo de componentes}

Cada agente pode atuar como cliente e servidor. A seguir, o diagrama de arquitetura desenhado nativamente em \LaTeX/TikZ.

\begin{figure}[h!]
\centering
\begin{tikzpicture}[
  font=\sffamily\small,
  agent/.style={draw=corBorder, fill=corAgent, rounded corners=10pt, inner sep=8pt, minimum width=9.2cm},
  box/.style={draw, fill=corKey, rounded corners=6pt, inner sep=4pt, text width=8.2cm, align=left, font=\sffamily\footnotesize},
  title/.style={font=\sffamily\bfseries\small, text=corTitulo},
  arr/.style={-{Stealth[length=3mm]}, thick, draw=corBorder},
  escrow/.style={draw=corVermelho, fill=corEscrow, rounded corners=10pt, inner sep=8pt, minimum width=8.6cm}
]

% Agente A
\node[agent] (A) at (0,0) {};
\node[title, anchor=north west] at (A.north west) {Agente A — Remetente};

\node[box, anchor=north west, xshift=0.4cm, yshift=-0.5cm] (chaveA) at (A.north west) {
  \textbf{Chaves do Agente A}\\
  Ed25519 (identidade/assinatura) $\cdot$ X25519 (acordo de chaves)\\
  Par efemero X25519 (por mensagem)
};

\node[box, anchor=north west, xshift=0.4cm, yshift=-2.6cm] (encA) at (A.north west) {
  \textcolor{corVerde}{\textbf{encrypt\_to():}} cifra payload AES-256-GCM
};

\node[box, anchor=north west, xshift=4.7cm, yshift=-2.6cm] (sigA) at (A.north west) {
  \textcolor{corAzul}{\textbf{assina Ed25519:}} cabeçalho + ambos os wraps
};

\node[box, anchor=north west, xshift=0.4cm, yshift=-4.4cm] (replayA) at (A.north west) {
  \textcolor{corLaranja}{\textbf{Registro de replay (SQLite):}} UUID da mensagem + expiração; rejeita reenvio
};

% Seta A -> B
\draw[arr] (A.east) -- ++(1.4cm,0) node[midway, above, font=\sffamily\footnotesize] {Envelope JSON v3};

% Agente B
\node[agent, anchor=west] (B) at ($(A.east)+(2.6cm,0)$) {};
\node[title, anchor=north west] at (B.north west) {Agente B — Destinatário};

\node[box, anchor=north west, xshift=0.4cm, yshift=-0.5cm] (chaveB) at (B.north west) {
  \textbf{Chaves do Agente B}\\
  Ed25519 (identidade) $\cdot$ X25519 (acordo de chaves)
};

\node[box, anchor=north west, xshift=0.4cm, yshift=-2.6cm] (decB) at (B.north west) {
  \textcolor{corAzul}{\textbf{decrypt\_message():}} unwrap K (chave própria)
};

\node[box, anchor=north west, xshift=4.7cm, yshift=-2.6cm] (verB) at (B.north west) {
  \textcolor{corAzul}{\textbf{verifica assinatura:}} Ed25519 do remetente A
};

\node[box, anchor=north west, xshift=0.4cm, yshift=-4.4cm] (authB) at (B.north west) {
  \textcolor{corRoxo}{\textbf{Autorização (negar por padrão):}} (remetente, kind, name) registrado?
};

\node[box, anchor=north west, xshift=0.4cm, yshift=-6.2cm] (execB) at (B.north west) {
  \textcolor{corVerde}{\textbf{dispatch() $\rightarrow$ handler:}} tool/skill executada localmente
};

\node[box, anchor=north west, xshift=4.7cm, yshift=-6.2cm] (respB) at (B.north west) {
  \textcolor{corVerde}{\textbf{Resposta cifrada:}} de volta para o Agente A
};

\draw[arr] (B.south) -- ++(0,-0.8cm);

% Escrow
\node[escrow, anchor=north, yshift=-0.2cm] (org) at ($(B.south)!0.5!(A.south)+(0,-2.2cm)$) {
  \textbf{Organização — Chave de Escrow}
};

\node[box, anchor=north west, xshift=0.4cm, yshift=-0.5cm] (orgKey) at (org.north west) {
  \textcolor{corVermelho}{\textbf{org\_private.pem (X25519):}} guarda-se somente com a empresa (HSM/ACL)
};

\node[box, anchor=north west, xshift=0.4cm, yshift=-1.8cm] (orgDec) at (org.north west) {
  \textcolor{corVermelho}{\textbf{org\_decrypt\_message():}} recupera K via wrap\_org $\rightarrow$ auditoria interna
};

\draw[arr, color=corVermelho] (B.south) -- ++(0,-1.6cm) -| ($(org.west)+(0.4cm,0)$);

\end{tikzpicture}
\caption{Figura 1 — Arquitetura do agente: remetente cifra, destinatário valida e executa, empresa descriptografa via escrow.}
\label{fig:arquitetura}
\end{figure}

\subsubsection*{4.3 Estrutura do envelope (v3)}

As mensagens são representadas em JSON e contêm o envelope abaixo. Os campos de roteamento (\texttt{v} \dots \texttt{capability}) são metadados visíveis; o conteúdo de negócios fica cifrado.

\begin{table}[h!]
\centering
\small
\begin{tabular}{p{3.4cm}p{2.2cm}p{10.4cm}}
\toprule
\rowcolor{corTitulo}\textcolor{white}{\textbf{Campo}} & \textcolor{white}{\textbf{Tipo}} & \textcolor{white}{\textbf{Descrição}} \\
\midrule
\texttt{v} & número & Versão do protocolo (3). \\
\texttt{msg\_id} & UUID & Identificador único da mensagem (controle de replay). \\
\texttt{org} & texto & Identificador da organização. \\
\texttt{sender} & texto & ID do remetente. \\
\texttt{recipient} & texto & ID do destinatário. \\
\texttt{issued} & inteiro & Instante de emissão (epoch segundos). \\
\texttt{expires} & inteiro & Prazo de validade (epoch segundos). \\
\texttt{sender\_x25519\_pub} & base64 & Chave pública X25519 efemera do remetente para esta mensagem. \\
\texttt{capability} & objeto & \{kind: tool|skill, name: …\} — roteamento, não secreto. \\
\texttt{nonce} & base64 & Nonce de 12 bytes do AES-256-GCM do payload. \\
\texttt{ciphertext} & base64 & Payload cifrado com AES-256-GCM. \\
\texttt{wrap\_recipient} & base64 & Chave simétrica K encapsulada para o destinatário. \\
\texttt{wrap\_org} & base64 & Chave simétrica K encapsulada para a empresa (escrow). \\
\texttt{signature} & base64 & Assinatura Ed25519 sobre o cabeçalho completo. \\
\bottomrule
\end{tabular}
\caption{Campos do envelope v3.}
\label{tab:envelope}
\end{table}

Payload interno (cifrado, JSON): \texttt{\{sender, recipient, capability, arguments\}}.

\subsubsection*{4.4 Sobre o Código}

Implementado em Python 3.11+. Dependência: biblioteca \texttt{cryptography} ($\geq$46, $<$47). A implementação usa a biblioteca padrão para JSON, SQLite, UUID e tempo, e a biblioteca \texttt{cryptography} para:
\begin{itemize}[leftmargin=*]
  \item X25519 (acordo de chaves), Ed25519 (assinaturas), AES-256-GCM (cifragem simétrica), HKDF-SHA256 (derivação).
  \item Geração de chaves (\texttt{scripts/generate\_chave\_keys.py}): um par de escrow X25519 por organização; um par Ed25519 + um par X25519 por agente.
  \item Cifragem (\texttt{cripto\_agent/chave\_protocol.py}): \texttt{encrypt\_message()}, \texttt{decrypt\_message()}, \texttt{org\_decrypt\_message()}.
  \item Agente (\texttt{SecureAgent}): registro de capabilities, catálogo filtrado por permissão, despacho com controle de replay persistente e descriptografia por escrow.
\end{itemize}

Não há carregamento de código remoto, instalação de skills via mensagem ou alteração de permissões por meio da rede. Cada handler recebe um objeto JSON, valida argumentos e retorna um resultado serializável em JSON.

% ================= 5. MATEMÁTICA =================
\sect{5. Estrutura Criptográfica Original — Matemática e Funcionamento}

O protocolo v3 adota uma estrutura de chaves original (estilo chaves), na qual a organização detém uma chave de descriptografia para segurança interna. Abaixo, explicamos a matemática por trás de cada primitivo.

\subsubsection*{5.1 Acordo de chaves — Curve25519 / X25519 (ECDH)}

Cada agente e a organização possuem um par de chaves X25519 sobre a curva elíptica Curve25519. Se um agente possui chave privada escalar $s$ e chave pública $P = s \cdot G$ (ponto base da curva), e outro possui par $(s', P')$, o segredo compartilhado é:

\begin{equation}
S = s \cdot P' \;=\; s' \cdot P \qquad \text{(pontos sobre Curve25519)}
\end{equation}

A dificuldade de deduzir $S$ a partir de $P$ e $P'$ (sem conhecer $s$ ou $s'$) baseia-se no problema do logaritmo discreto em Curve25519, considerado computacionalmente intratável. Esse segredo $S$ nunca é usado diretamente como chave de cifragem; ele passa por uma derivação de chaves.

\subsubsection*{5.2 Derivação de chaves — HKDF-SHA256}

A partir do segredo compartilhado $S$, derivamos chaves simétricas com HKDF-SHA256, que tem duas etapas:

\begin{align}
\mathrm{PRK} &= \mathrm{HMAC\text{-}SHA256}(\text{salt}, \mathrm{IKM}) && \text{(extração)}\\
\mathrm{OKM} &= \mathrm{HMAC\text{-}SHA256}(\mathrm{PRK},\; \mathrm{info}\!\parallel\!0x01) && \text{(expansão, primeiro bloco)}
\end{align}

No protocolo, usamos $\mathrm{IKM} = S$, salt vazio e contextos ($\mathrm{info}$) distintos:

\begin{align}
K_{\mathrm{rec}} &= \mathrm{HKDF}(S_{\mathrm{rec}},\; \text{\texttt{org-chave/recipient/v3}}) && \text{encapsulamento para o destinatário}\\
K_{\mathrm{org}} &= \mathrm{HKDF}(S_{\mathrm{org}},\; \text{\texttt{org-chave/escrow/v3}}) && \text{encapsulamento para a empresa}
\end{align}

Assim, a mesma mensagem produz duas chaves simétricas independentes derivadas de dois acordos ECDH diferentes: um com a chave pública do destinatário e outro com a chave pública de escrow da organização.

\subsubsection*{5.3 Cifragem do payload — AES-256-GCM}

Cada mensagem gera uma chave simétrica aleatória $K$ (256 bits) e um nonce de 96 bits (12 bytes). O payload $P$ é cifrado como:

\begin{equation}
C = \mathrm{AES\text{-}256\text{-}GCM\text{-}Encrypt}(K,\; \text{nonce},\; P,\; \mathrm{AAD})
\end{equation}

Onde $\mathrm{AAD}$ (\emph{Authenticated Additional Data}) é o cabeçalho de roteamento canônico, de modo que qualquer alteração nos campos de cabeçalho invalida a autenticação do cifrado. O AES-GCM fornece confidencialidade e integridade em uma só operação, produzindo o texto cifrado $C$ e uma etiqueta de autenticação (tag) de 128 bits.

\subsubsection*{5.4 Encapsulamento duplo da chave $K$}

A chave $K$ é cifrada duas vezes, cada uma com uma chave derivada de um acordo ECDH diferente:

\begin{align}
\mathrm{wrap}_{\mathrm{recipient}} &= \mathrm{AES\text{-}256\text{-}GCM}(K_{\mathrm{rec}},\; \text{nonce}_1,\; K,\; \mathrm{aad}_{\mathrm{recipient}})\\
\mathrm{wrap}_{\mathrm{org}} &= \mathrm{AES\text{-}256\text{-}GCM}(K_{\mathrm{org}},\; \text{nonce}_2,\; K,\; \mathrm{aad}_{\mathrm{org}})
\end{align}

Apenas o destinatário (possuidor de $B_{\mathrm{priv}}$, que recompõe $S_{\mathrm{rec}} = B_{\mathrm{priv}} \cdot E_{\mathrm{pub}} \rightarrow K_{\mathrm{rec}}$) e a empresa (possuidora de $\mathrm{Org}_{\mathrm{priv}}$, que recompõe $S_{\mathrm{org}} = \mathrm{Org}_{\mathrm{priv}} \cdot E_{\mathrm{pub}} \rightarrow K_{\mathrm{org}}$) conseguem recuperar $K$ e, portanto, descifrar o payload. O encapsulamento vincula o par efemero $E_{\mathrm{pub}}$ enviado no cabeçalho.

\subsubsection*{5.5 Assinatura — Ed25519}

O remetente assina todo o cabeçalho (incluindo $\mathrm{wrap}_{\mathrm{recipient}}$ e $\mathrm{wrap}_{\mathrm{org}}$) com sua chave privada Ed25519:

\begin{equation}
\sigma = \mathrm{Sign}(A_{\mathrm{priv\_ed}},\; \text{cabeçalho canônico})
\end{equation}

A verificação usa a chave pública Ed25519 cadastrada do remetente:

\begin{equation}
\mathrm{Verify}(A_{\mathrm{pub\_ed}},\; \sigma,\; \text{cabeçalho canônico}) = \text{válido / inválido}
\end{equation}

Isso garante autenticidade, integridade do cabeçalho e não repúdio: o remetente não pode negar ter enviado a mensagem, e qualquer alteração nos encapsulamentos de chave invalida a assinatura antes mesmo da descifragem.

\subsubsection*{5.6 Escrow organizacional — chave de descriptografia interna da empresa}

A organização possui um par X25519 de escrow. A chave privada $\mathrm{Org}_{\mathrm{priv}}$ fica sob guarda exclusiva da empresa (idealmente HSM), enquanto a pública $\mathrm{Org}_{\mathrm{pub}}$ é distribuída a todos os agentes. Toda mensagem é encapsulada para $\mathrm{Org}_{\mathrm{pub}}$ ($\mathrm{wrap}_{\mathrm{org}}$), de modo que a empresa pode, a qualquer momento:

\begin{enumerate}[leftmargin=*]
  \item Recuperar $K = \mathrm{Unwrap}(\mathrm{Org}_{\mathrm{priv}},\; \mathrm{wrap}_{\mathrm{org}},\; E_{\mathrm{pub}})$ usando o par efemero $E_{\mathrm{pub}}$ do cabeçalho;
  \item Descifrar o payload com $K$ e verificar sua integridade (etiqueta AES-GCM);
  \item Auditar o conteúdo da mensagem sem quebrar o modelo agente-a-agente.
\end{enumerate}

O destinatário ainda consegue perceber quando a empresa precisou intervir (campo \texttt{decrypted\_via\_org}), e a empresa não pode falsificar mensagens, pois não possui a chave privada Ed25519 dos agentes (apenas valida assinaturas alheias).

\begin{figure}[h!]
\centering
\begin{tikzpicture}[
  font=\sffamily\small,
  colA/.style={draw=corBorder, fill=corAgent, rounded corners=10pt, inner sep=8pt, minimum width=9.6cm},
  box/.style={draw, fill=corKey, rounded corners=6pt, inner sep=4pt, text width=9.0cm, align=left, font=\sffamily\footnotesize},
  title/.style={font=\sffamily\bfseries\small, text=corTitulo},
  arr/.style={-{Stealth[length=3mm]}, thick, draw=corBorder},
  enc/.style={draw=corAzul, fill=corDestaque, rounded corners=8pt, inner sep=6pt, text width=9.0cm, align=left, font=\sffamily\footnotesize}
]

\node[colA] (R) at (0,0) {};
\node[title, anchor=north west] at (R.north west) {Remetente (Agente A)};

\node[box, anchor=north west, xshift=0.4cm, yshift=-0.5cm] (p) at (R.north west) {
  \textbf{Payload (negócios)}\\
  \texttt{\{sender, recipient, capability, arguments\}}
};

\node[box, anchor=north west, xshift=0.4cm, yshift=-1.9cm] (eph) at (R.north west) {
  \textbf{1. Par efemero X25519 (E)}\\
  $E_{\mathrm{priv}}$, $E_{\mathrm{pub}}$ gerados aleatoriamente por mensagem
};

\node[box, anchor=north west, xshift=0.4cm, yshift=-3.3cm] (sym) at (R.north west) {
  \textbf{2. Chave $K$ $\leftarrow$ AES-256 aleatória (256 bits)}\\
  $C = \mathrm{AES\text{-}256\text{-}GCM}(K,\; \text{nonce},\; \text{payload},\; \mathrm{AAD}=\text{cabeçalho})$
};

\node[enc, anchor=north west, xshift=0.4cm, yshift=-4.9cm] (hkdf) at (R.north west) {
  \textbf{3. Encapsulamento duplo de $K$ (ECDH + HKDF-SHA256)}\\
  $S_{\mathrm{rec}} = E_{\mathrm{priv}} \cdot B_{\mathrm{pub}} \;\rightarrow\; K_{\mathrm{rec}} = \mathrm{HKDF}(S_{\mathrm{rec}},\; \text{\texttt{org-chave/recipient/v3}})$\\
  $S_{\mathrm{org}} = E_{\mathrm{priv}} \cdot \mathrm{Org}_{\mathrm{pub}} \;\rightarrow\; K_{\mathrm{org}} = \mathrm{HKDF}(S_{\mathrm{org}},\; \text{\texttt{org-chave/escrow/v3}})$\\
  $\mathrm{wrap}_{\mathrm{recipient}} = \mathrm{AES\text{-}GCM}(K_{\mathrm{rec}},\; \text{nonce}_1,\; K)$\\
  $\mathrm{wrap}_{\mathrm{org}} = \mathrm{AES\text{-}GCM}(K_{\mathrm{org}},\; \text{nonce}_2,\; K)$
};

\node[box, anchor=north west, xshift=0.4cm, yshift=-8.6cm] (sig) at (R.north west) {
  \textbf{4. Assinatura Ed25519}\\
  $\sigma = \mathrm{Sign}(A_{\mathrm{priv\_ed}},\; \text{cabeçalho} + \mathrm{wrap}_{\mathrm{recipient}} + \mathrm{wrap}_{\mathrm{org}})$\\
  Envelope = $\{\texttt{v, msg\_id, org, sender, recipient, issued, expires,}$\\
  \phantom{Envelope = }$\texttt{E\_pub, capability, nonce, ciphertext, wrap\_recipient, wrap\_org, signature}\}$
};

\draw[arr] (R.east) -- ++(1.5cm,0) node[midway, above, font=\sffamily\footnotesize] {rede organizacional (JSON)};

\node[colA, anchor=west] (D) at ($(R.east)+(2.7cm,0)$) {};
\node[title, anchor=north west] at (D.north west) {Destinatário (Agente B)};

\node[box, anchor=north west, xshift=0.4cm, yshift=-0.5cm] (v1) at (D.north west) {
  \textbf{1. Verifica assinatura $\sigma$}\\
  $\mathrm{Verify}(A_{\mathrm{pub\_ed}},\; \sigma,\; \text{cabeçalho}) \rightarrow$ autenticidade / não repúdio
};

\node[box, anchor=north west, xshift=0.4cm, yshift=-2.1cm] (v2) at (D.north west) {
  \textbf{2. Validade temporal / organização}\\
  \texttt{issued/expires} dentro da janela (com tolerância de relógio); \texttt{org} confere
};

\node[box, anchor=north west, xshift=0.4cm, yshift=-3.7cm] (v3) at (D.north west) {
  \textbf{3. Recupera $K$ (unwrap)}\\
  tenta: $K = \mathrm{Unwrap}(B_{\mathrm{priv}},\; \mathrm{wrap}_{\mathrm{recipient}},\; E_{\mathrm{pub}})$\\
  se falhar e houver $\mathrm{Org}_{\mathrm{priv}}$: $K = \mathrm{Unwrap}(\mathrm{Org}_{\mathrm{priv}},\; \mathrm{wrap}_{\mathrm{org}},\; E_{\mathrm{pub}})$ $\leftarrow$ via escrow\\
  se ambos falharem $\rightarrow$ rejeita (\texttt{cannot\_unwrap})
};

\node[box, anchor=north west, xshift=0.4cm, yshift=-6.3cm] (v4) at (D.north west) {
  \textbf{4. Descifra payload}\\
  $\text{payload} = \mathrm{AES\text{-}GCM\text{-}decrypt}(K,\; \text{nonce},\; C,\; \mathrm{AAD}=\text{cabeçalho})$\\
  confere sender/recipient/capability contra o cabeçalho
};

\node[box, anchor=north west, xshift=0.4cm, yshift=-7.9cm] (v5) at (D.north west) {
  \textbf{5. Replay + autorização + execução}\\
  UUID persistido (SQLite) $\rightarrow$ rejeita repetição\\
  (remetente, kind, name) com permissão explícita $\rightarrow$ handler executado
};

\node[box, anchor=north west, xshift=0.4cm, yshift=-9.7cm] (v6) at (D.north west) {
  \textbf{6. Resposta para A}\\
  resultado cifrado; empresa pode descriptografar via escrow
};

\end{tikzpicture}
\caption{Figura 2 — Fluxo de cifragem de uma mensagem v3 (remetente $\rightarrow$ destinatário).}
\label{fig:fluxo}
\end{figure}

% ================= 6. MANUAL =================
\sect{6. Manual de Utilização para Usuários Contemplados}

O manual segue a estrutura de guia de instruções por tarefa, com exceções e problemas potenciais, conforme o template.

\subsubsection*{6.1 Tarefa A — Gerar chaves da organização e dos agentes}

\begin{enumerate}[leftmargin=*]
  \item Acesse a pasta do projeto no Windows e rode:
    \begin{lstlisting}
python scripts\generate_chave_keys.py --out .local\chave
    \end{lstlisting}
  \item O gerador cria \texttt{org\_private.pem} (somente a empresa), \texttt{org\_public.pem} e um par de chaves (Ed25519 + X25519) por agente, além de \texttt{config.json}.
\end{enumerate}

\textbf{Exceções ou potenciais problemas:} se a pasta de saída já existe, o gerador recusa sobrescrever (evita perda de chaves). Use outro \texttt{--out} ou remova a pasta.

Guarde \texttt{org\_private.pem} apenas com a empresa (escrow). Não compartilhe chaves privadas.

\subsubsection*{6.2 Tarefa B — Executar a demonstração de ponta a ponta}

\begin{enumerate}[leftmargin=*]
  \item Rode o exemplo executável:
    \begin{lstlisting}
python examples\chave_demo.py
    \end{lstlisting}
  \item O agente A cifra uma mensagem para o agente B; B descifra com sua própria chave, executa a tool/skill e responde. A empresa descriptografa a mesma mensagem via escrow (\texttt{org\_decrypt\_message}) e o replay da mesma mensagem é rejeitado.
\end{enumerate}

\textbf{Exceções:} assinatura inválida $\rightarrow$ \texttt{bad\_signature}; organização diferente $\rightarrow$ \texttt{org\_mismatch}; chave não recuperável $\rightarrow$ \texttt{cannot\_unwrap}; UUID repetido $\rightarrow$ \texttt{replay}.

\subsubsection*{6.3 Tarefa C — Testes automatizados}

\begin{enumerate}[leftmargin=*]
  \item Rode a suíte de testes (cobre o protocolo v1 HTTPS/mTLS e o v3 Chave-Organização):
    \begin{lstlisting}
python -m unittest discover -s tests -v
    \end{lstlisting}
  \item Esperado: todos os testes passam, validando round-trip de cifragem, escrow da empresa, assinatura, expiração, organização, replay e autorização.
\end{enumerate}

% ================= 7. CONSIDERAÇÕES =================
\sect{7. Considerações Finais}

O protocolo Chave-Organização v3 entrega uma estrutura criptográfica original, com chave de descriptografia interna da empresa (escrow), confidencialidade ponta-a-ponta (AES-256-GCM), autenticidade e não repúdio (Ed25519), acordo de chaves (X25519/ECDH + HKDF-SHA256) e controle de replay persistente. É um protótipo funcional que requer homologação entre máquinas reais e recursos operacionais (gerência centralizada de identidades, rotação/revogação de chaves, HSM para a chave de escrow, isolamento de handlers e testes de rede distribuída) antes de adoção em produção.

Este relatório segue o Modelo Estrutural para Relatório Final do Projeto, com nome do projeto, breve descrição, quatro cenários de visão de projeto, documentação técnica (requisitos, arquitetura, envelope, código) e manual de utilização, incluindo diagramas nativos em TikZ e a explicação matemática dos primitivos criptográficos.

\end{document}
